% Figure: on-axis field of a single coil, of two coils spaced by their
% radius (the Helmholtz condition), and of two coils spaced too far apart.
% The Helmholtz pair has a flat plateau at the centre because the second
% axial derivative of the summed field vanishes there.
\begin{figure}[htbp]
\centering
\begin{tikzpicture}[>=Latex]
\begin{axis}[
  width=12cm, height=7cm,
  xlabel={axial position $x/a$ (coil radius $a$)},
  ylabel={on-axis field $B(x)/B_0$ (centre value)},
  xmin=-1.5, xmax=1.5,
  ymin=0, ymax=1.15,
  legend pos=south west,
  legend cell align=left,
  font=\scriptsize,
  grid=both,
  grid style={enggray!20},
]
% single coil of radius 1 at origin: B ~ 1/(1+x^2)^{3/2}, normalised to 1 at x=0
\addplot[engred, line width=1.1pt, domain=-1.5:1.5, samples=120]
  {1/(1+x^2)^(1.5)};
\addlegendentry{single coil}
% Helmholtz pair at x = +-1/2, normalised so centre = 1
% raw sum s(x) = 1/(1+(x-0.5)^2)^{1.5} + 1/(1+(x+0.5)^2)^{1.5}
% centre value s(0) = 2/(1.25)^{1.5} = 2/1.3975 = 1.4311
\addplot[engblue, line width=1.4pt, domain=-1.5:1.5, samples=160]
  {(1/(1+(x-0.5)^2)^(1.5) + 1/(1+(x+0.5)^2)^(1.5))/1.43108};
\addlegendentry{Helmholtz pair, spacing $=a$}
% too-far pair at x = +-1.0, normalised so centre = 1
% s(0) = 2/(2)^{1.5} = 2/2.8284 = 0.7071
\addplot[engamber, line width=1.1pt, dashed, domain=-1.5:1.5, samples=160]
  {(1/(1+(x-1.0)^2)^(1.5) + 1/(1+(x+1.0)^2)^(1.5))/0.70711};
\addlegendentry{pair, spacing $=2a$ (too far)}
% one-percent band around unity
\addplot[enggray, line width=0.4pt, domain=-1.5:1.5, samples=2, forget plot]
  {0.99};
\addplot[enggray, line width=0.4pt, domain=-1.5:1.5, samples=2, forget plot]
  {1.01};
\node[enggray, anchor=south west, font=\scriptsize] at (axis cs:-1.45,1.012)
  {$\pm 1\%$ band};
\end{axis}
\end{tikzpicture}
\caption[Helmholtz-coil field uniformity]{On-axis field of a single coil, of
a Helmholtz pair spaced by exactly one coil radius, and of a pair spaced too
far apart, each normalised to its own centre value. The single coil peaks
sharply at its plane and falls off on either side. Two coils spaced by their
radius produce a broad flat plateau at the midpoint, because the second axial
derivative of the summed field vanishes there when the spacing equals the
radius, leaving the field uniform to better than one percent over a central
region roughly a third of the radius wide. Spacing the coils too far apart
splits the plateau into two humps with a dip at the centre, the sign that the
uniformity condition has been overshot. The Helmholtz geometry is the working
tool for a region of controlled, uniform field on a bench, used to cancel the
Earth's field over a sample or to calibrate a magnetometer.}
\label{fig:vol04ch08:helmholtz}
\end{figure}
